%Dem jeg ved jeg bruger%%%%%%%%%%%%%%%%%%
\makeatletter
\usepackage{geometry}
\usepackage{graphicx}
\usepackage{titling}
\usepackage{lastpage}
\usepackage{fancyhdr}
\usepackage[T1]{fontenc}
% \usepackage[danish]{babel} %dansk orddeling og overskrifter
\usepackage{indentfirst} %indents i starten af paragrapher
\geometry{
  a4paper,
  total={170mm,257mm},
 left=20mm,
 top=20mm,
 }
 %\usepackage{apacite}%citations
 \usepackage{url}
 %%%%%%%%%%%%%%%%%%%%%%%%%%%%%%%%%%%%%%%
 \newcommand{\sd}[1]{\textsubscript{#1}}
 \newcommand{\su}[1]{\textsuperscript{#1}}
 \usepackage[makeroom]{cancel}
\usepackage{pgfplots, pgfplotstable}%plots
\pgfplotsset{
axis lines = left,
width=0.8\linewidth,
height =0.4\textwidth,
compat = newest,
grid = both,
grid style={line width=.1pt, draw=gray!10},
minor tick num=5,
major grid style={line width=.2pt,draw=gray!50},
legend style={cells={align=center}}
}%define standert plot peramiters
\pgfkeys{
     /pgf/number format/.cd,% <- changes the prefix for the following options
     1000 sep={}
  }
\usepackage{plimsoll}%stander state symbol
\usetikzlibrary{fpu}
\usepackage{fp}
\pgfkeys{/pgf/fpu/output format=float}
%siljas magi%%%%%%%%%%%%%%
\usepackage{amsmath}

\usepackage{mathtools} %ligninger og den slags
\usepackage{graphicx} %indsætning af billeder og indstillinger til disse
\usepackage{amssymb} %stor pakke af symboler til matematikken
\usepackage{setspace} %indstilling af linjeafstand
%\usepackage[left=3cm,right=3cm,top=3.5cm,bottom=3cm]{geometry} %indstilling af margener
%\usepackage[table]{xcolor} %mulighed for farvede felter i skemaer
\usepackage{float}
\usepackage{listings}
\usepackage{multirow}
\usepackage{circuitikz}
\usepackage{caption}
\usepackage{textcomp}\usepackage{titlesec}
\usetikzlibrary{arrows,shapes.gates.logic.US,shapes.gates.logic.IEC,calc}
\usetikzlibrary{circuits.logic.US} % TiKZ Library for US Logic Circuits.
\tikzstyle{branch}=[fill,shape=circle,minimum size=3pt,inner sep=0pt]
\lstset{language=VHDL}
\lstset{basicstyle=\footnotesize\ttfamily,breaklines=true,numbers=none}
\PassOptionsToPackage{hyphens}{url}\usepackage{hyperref}
\usepackage{pdfpages}
\usepackage[justification=centering]{caption}
\usepackage{subcaption}
\usepackage{wrapfig}
\usepackage{multicol}
\usepackage{comment}
\usepackage[activate={true,nocompatibility},final,tracking=true,kerning=true,spacing=true,factor=1100,stretch=10,shrink=10]{microtype}
\usepackage[space]{grffile}
% \newcommand{\Mypm}{\mathbin{\tikz [x=1.4ex,y=1.4ex,line width=.1ex] \draw (0.0,0) -- (1.0,0) (0.5,0.08) -- (0.5,0.92) (0.0,0.5) -- (1.0,0.5);}}%
% \setlength{\parindent}{0em}
\usepackage{listings}
\usepackage[most]{tcolorbox}

\usepackage[minimal = false,formula = {chemformula}]{chemmacros}
\usepackage{chemformula}
\usepgfplotslibrary{patchplots}
\tikzset{
  ctrlpoint/.style={%
    draw=black,
    circle,
    inner sep=0.1,
    minimum width=0.4ex,
    fill=black
  }
}
\DeclareRobustCommand{\vol}{\text{\volumedash}V}
\newcommand{\volumedash}{%
  \makebox[0pt][l]{%
    \ooalign{\hfil\hphantom{$\m@th V$}\hfil\cr\kern0.08em--\hfil\cr}%
  }%
}
\makeatother
\usepackage{placeins}